\documentclass[]{article} 
\usepackage[english,catalan]{babel}
\usepackage{csquotes}
\usepackage{graphicx} % Imatges
\usepackage{fancyhdr}
\usepackage{ragged2e}
\usepackage{hyperref}
\usepackage{textcomp}
\usepackage[x11names, rgb]{xcolor}
\usepackage{listings}
\usepackage{enumitem}
\usepackage{pdfpages}
\usepackage[notes,backend=biber,defernumbers=true]{biblatex-chicago}
\bibliography{MemoriaPractiques}
\usepackage{setspace}
\usepackage{rotating}
\usepackage{censor}
\usepackage[paper=portrait,pagesize]{typearea}
\usepackage{pdflscape}
\usepackage{geometry}
    \geometry{
    a4paper,
    inner=3.75cm,
    outer=1.75cm,
    top=3cm,
    bottom=3.5cm,
    headheight=1.5cm,
    footskip=42pt}

\hypersetup{
    colorlinks,
    linkcolor={blue!50!black},
    citecolor={blue!50!black},
    urlcolor={blue!80!black}
}
   

\title{KAIZEN STRATEGIES FOR IMPROVING PLANNING AND PRODUCTION FLOW IN PHARMACEUTICAL ENVIRONMENTS}
\author{Miquel Borrell I Candelich}
\date{August 2023}


\setlength{\parindent}{0pt}

\lstset{
    % How/what to match
    sensitive=true,
    % Border (above and below)
    frame=lines,
    % Extra margin on line (align with paragraph)
    xleftmargin=\parindent,
    % Put extra space under caption
    belowcaptionskip=1\baselineskip,
    % Colors
    backgroundcolor=\color{sbase3},
    basicstyle=\color{sbase00}\ttfamily\scriptsize,
    keywordstyle=\color{scyan},
    commentstyle=\color{sbase1},
    stringstyle=\color{sblue},
    numberstyle=\color{sred},
    identifierstyle=\color{sorange},
    % Break long lines into multiple lines?
    breaklines=true,
    % Show a character for spaces?
    showstringspaces=false,
    tabsize=2
}

\usepackage[T1]{fontenc}
\usepackage[normalem]{ulem}

\newif\ifconfid
\confidtrue                   % comment this line out to unveil all hidden text
\colorlet{confidcolor}{black} % change the colour of boxes here (e.g. "red")


% ----------- implementation details -----------
\makeatletter

% from Martin's Scharrer's answer (https://tex.stackexchange.com/a/16004/21891)
\def\confid@UL@putbox{%
  \ifx\UL@start\@empty%
  \else % not inner
    \vrule\@width\z@ \LA@penalty\@M
    {\UL@skip\wd\UL@box \UL@leaders \kern-\UL@skip}%
    \phantom{\box\UL@box}%
  \fi
}

\newcommand\confidentiel{}%

\ifconfid
\renewcommand\confidentiel[1][confidcolor]{%
 \bgroup%
 \let\UL@putbox\confid@UL@putbox%
 \markoverwith{\hbox to.01em{\hss\textcolor{#1}{|}\hss}}\ULon%
}
\fi

\lst@AddToHook{Init}{%
  \ifconfid%
    \lstset{moredelim=[is][\confidentiel]``,keepspaces}%
  \else
    \lstset{moredelim=**[is][]``}%
  \fi
}
\makeatother
% ----------- end of implementation details -----------

\begin{document}
\onehalfspacing
%Peus de pàgina
\pagestyle{fancy}
\fancyfoot[R]{\includegraphics[]{footer4.png}}
\fancyfoot[C]{\thepage}
\fancyhead[R]{Kaizen Strategies For Improving Planning And Production Flow In Pharmaceutical Environments}
\fancyhead[L]{}



\begin{titlepage}

\includegraphics[width=0.4\textwidth]{LogoUPC.png}
\begin{center}
       \vspace*{1cm}

       \textbf{Bachelor's Thesis}

       \textbf{KAIZEN STRATEGIES FOR IMPROVING PLANNING AND PRODUCTION FLOW IN PHARMACEUTICAL ENVIRONMENTS}

       \vspace{0.5cm}
        \textbf{Double degree in Chemical Engineering and Industrial Electronics and Automation.}\\
            
       \vspace{1.5cm}

       \textbf{In colaboration with: Kaizen Institute}\\
       Location: Barcelona\\

\end{center}
\includegraphics[]{BannerEEBE.png}
\begin{center}       
\textbf{Report and Annex}\\
\end{center}
       \vfill
\textbf{Author}:Miquel Borrell i Candelich\\
\textbf{EEBE Supervisor:} Moisés Graells\\
\textbf{EEBE Co-Supervisor:} Roque Lopez Paricio\\
\textbf{Collaborating entity tutor:} Sergi Bernà\\
\textbf{Call:} June, 2025\\
            
       \vspace{0.8cm}
            
\end{titlepage}



\newpage
\newpage

\justifying
\section{Abstract}\label{abstract}
\newpage

\section{Resum}\label{resum}
\newpage

\section{Resumen}\label{resumen}
\newpage

\section{Glossary}
\newpage

% TABLE OF CONTENTS
\tableofcontents

\section{Preface}\label{preface}
\subsection{The KAIZEN Philosofy}\label{philosofy}
After World War II, Toyoda Kiichiro, then president of the Toyota Motor Company, saw the need
to catch up with America to assure the survival of the automobile industry of Japan. Taiichi Ohno,
Toyota’s engineer, had a challenge: how to build a diversified line of products with the limited
equipment Toyota possessed at the time. Instead of using batch production, he designed a system
of integrated production that used one-piece flow, demonstrating that a trade-off between quality
and productivity does not need to exist (Netland and Powell, 2016). Aligned with other practices,
Toyota Production System (TPS) was born to defy the premise that bigger batches mean lower
costs.

TPS, which ultimate goal is the absolute elimination of waste, is based in two pillars: just-in-
time (JIT) and autonomation (automation with a human touch). The lack of the necessary capital
resources to support high inventory levels was solved with the implementation of a just-in-time
pull system. JIT means "to produce the necessary units at the necessary quantities at the necessary
time". A company establishing this flow throughout can approach zero inventory. With autonoma-
tion, Toyota intended to create a built-in automatic checking system against small abnormalities.
With this prevention mechanism, defective products were not produced. That was not the only
advantage. While the machine is working properly, employees are liberated. Consequently, one
worker could attend multiple machines, increasing production efficiency. Furthermore, autonoma-
tion came as a facilitator of improvement. Stopping the machine creates awareness and, through
the total comprehension of the problem, improvement measures can be taken. (Ohno, 1988)
Lean, based on TPS, is a methodology that accentuates customer needs, to improve quality
and reduce costs through continuous improvement and employee involvement (Graban, 2008).
Kaizen, a Japanese word, refers to the concept of continuous improvement. It does not intend
to be mistaken as innovation or disruption and elevated costs. On the other hand, Kaizen is about
small and subtle improvements. It is a low-cost low-risk process that assures incremental progress
and sustainable changes in the long term.

At the beginning of the 21st century, Toyota Motor Company surpassed General Motors and
became the world leader in automobile production, expanding and elevating the Kaizen concept
for its key role in the company’s success. Kaizen’s methodology is based on five
fundamental principles:
\begin{itemize}
  \item Create customer value
    \subitem Value is what the customer is willing to pay for. Using a market-in approach to make
informed decisions based on what the customers want and offer them that in the best way
possible, improving customer experience.
  
  \item Create flow efficiency
\end{itemize}

\clearpage
\newpage
%\printbibliography[heading=bibnumbered]
\end{document}
